\section{Cyclic relations on the leading free-word errors}
\label{sec:cyclic-obstructions}

The space \(\mathcal G\) controls the linear corrections available
at several blocks. We must show that the errors to be corrected
satisfy its relations for individual cyclic orbits of free words.
Commuting scalar substitutions would combine too many words. We
instead substitute matrices, repeat each block, and attach hubs
whose derivative columns identify the possible left nullspace.

Traces of words and their cyclic symmetry are classical in matrix
invariant theory; see \cite[\S4]{Procesi1976}. Here the response
relations must first be proved to annihilate every matrix trace test.
An explicit directed-cycle calculation then isolates each cyclic orbit.

Fix \(m\geq1\) independent generic size-\(b\) points \(P_i\).
Write \(I_i\) for their selected root labels, so \(|I_i|=a\), and
let \(\mathcal G\subseteq\bigoplus_i H_{P_i}^*\) be the space
of Proposition~\ref{prop:compatible-relations}. As before, all
scalars may be extended to a common field containing the required
splitting data.

\begin{theorem}[Cyclic compatibility of the error]
\label{thm:cyclic-compatibility}
Let \(Z(s_1,s_2)\) be a pair of formal series in two free letters,
with constant term \(\operatorname{diag}(P_1,\ldots,P_m)\).
Suppose that the fixed-root Schur expressions \(S_{i,h}\), for
every block \(i\) and every root label \(h\), vanish in every
degree below \(\ell\), where \(\ell\geq1\).
By transfer, their degree-\(\ell\) coefficients are scalar on
their root spaces; write
\[
 [s^\omega]S_{i,h}=e_{i,h,\omega}\pi_{i,h},
 \qquad |\omega|=\ell.
\]
For each word the scalar tuple at block \(i\) belongs to its
active space. For every cyclic orbit \(\mathcal O\) of words of
length \(\ell\) and every \(c\in\mathcal G\),
\begin{equation}
 \sum_{\omega\in\mathcal O}
 \sum_{i=1}^m\sum_{h\in I_i}
 c_{i,h}e_{i,h,\omega}=0.
 \label{eq:cyclic-error-relation}
\end{equation}
The orbit sum runs over distinct words, including when a word has
a nontrivial rotational stabilizer.
\end{theorem}

We prove the theorem in three steps. First, arbitrary matrix
substitutions turn the free-word errors into matrices on a repeat
factor. Second, a derivative test shows that every relation in
\(\mathcal G\) annihilates their traces. Finally, directed-cycle
matrices isolate each cyclic orbit.

\begin{lemma}[Splitting repeated blocks]
\label{lem:repeated-splitting}
Fix \(L\geq1\) and arbitrary scalar matrices
\(S_1,S_2\in M_L(E)\). Substitute \(s_c=\delta S_c\) into
\(Z\), with a commuting Taylor parameter \(\delta\), and add
an order-\(\delta^\ell\) input loop \(Y_i\) in each old block
\(i\). Here \(Y_i\) is an arbitrary pair in
\(M_b(E)\otimes M_L(E)\). Put
\[
 B_{i,h}=\sum_{|\omega|=\ell}
 e_{i,h,\omega}S^\omega+D_{i,h}(Y_i),
\]
where \(D_{i,h}\) is applied entrywise to the repeat-factor
coefficients. One can choose the \(Y_i\) so that every \(B_{i,h}\),
for all root labels \(h\), has simple spectrum. For such a choice,
after a finite splitting extension the leaf output has branches
\[
 \lambda_{i,h,j}(\delta)
 =\lambda_{i,h}+\delta^\ell\theta_{i,h,j}
    +O(\delta^{\ell+1}),\qquad 1\leq j\leq L,
\]
each scalar of multiplicity \(t\), where \(\theta_{i,h,j}\)
are the eigenvalues of \(B_{i,h}\). Their spectral projections
are integral in \(\delta\) and reduce to
\(\pi_{i,h}\otimes\Pi_{i,h,j}\), where
\(\Pi_{i,h,j}\) are the rank-one spectral projections of
\(B_{i,h}\).
\end{lemma}

\begin{proof}
Matrix substitution is legitimate coefficientwise in \(\delta\).
The leading Schur expression on the repeated \((i,h)\)-space is
\(I_t\otimes B_{i,h}\): the original errors give the first sum,
and an order-\(\delta^\ell\) input loop contributes its first
response. All lower orders are zero.

For each label \(h\), \(D_{i,h}\) is a nonzero linear functional
on the input directions. Its matrix-valued extension is therefore
onto \(M_L(E)\). The discriminant of \(B_{i,h}\), as a polynomial
in the loop coordinates, is consequently not identically zero.
There are finitely many labels, so the product of these discriminants
is nonzero. Over the infinite field in use, or a scalar extension,
choose loops at which their product is nonzero. This choice is made
for all labels, not merely the selected ones. For \(L=1\) the
simple-spectrum condition is automatic.

Separate the old distinct roots by a formal conjugation equal to
the identity at \(\delta=0\). The separated \((i,h)\)-block is
\begin{equation}
 \lambda_{i,h}I+\delta^\ell(I_t\otimes B_{i,h})
 +O(\delta^{\ell+1}).
 \label{eq:old-cluster-expansion}
\end{equation}
Here is an exact comparison with the fixed-root Schur expression.
Write \(M=W-\lambda_{i,h}I\) in the fixed root/complement
decomposition, with blocks \(M_{uv}\). Its complementary block
\(M_{22}\) is invertible. If \(F\) is the separated root block
minus \(\lambda_{i,h}I\), its invariant columns can be written
\(\binom QR\), with \(Q=I+O(\delta)\), \(R=O(\delta)\), and
\(M\binom QR=\binom QR F\). Eliminating the lower row gives
\[
 \bigl(M_{11}-M_{12}M_{22}^{-1}M_{21}\bigr)Q
   =\bigl(Q-M_{12}M_{22}^{-1}R\bigr)F.
\]
The left parenthesis is the fixed-root Schur expression. Both
outer factors relating it to \(F\) reduce to identity at zero.
Consequently it and \(F\) have the same first nonzero coefficient.
Its vanishing below degree \(\ell\) forces the same vanishing
for \(F\), and its degree-\(\ell\) coefficient is
\(I_t\otimes B_{i,h}\). This proves
\eqref{eq:old-cluster-expansion}.

Within this block divide the change from \(\lambda_{i,h}I\) by
\(\delta^\ell\). Its constant term is \(I_t\otimes B_{i,h}\),
with \(L\) distinct eigenvalues. A second formal conjugation,
integral in \(\delta\), separates these slopes into blocks of
size \(t\). Across all old labels there are exactly
\[
 mLe_b=e_{bmL}
\]
such blocks. Their characteristic polynomials are pairwise coprime
over the Laurent fraction field: different old roots are already
distinct at order zero, and distinct slopes separate the blocks
inside one old root. The minimal-degree bound at size \(bmL\)
is \(e_{bmL}\). Each separated block contributes at least one
to that degree, so each contributes exactly one and is scalar.
The two integral conjugations also give integral spectral
projections with the stated limits.
\end{proof}

Keep a choice as in Lemma~\ref{lem:repeated-splitting}. Adjoin
\(L+1\) independent generic size-\(b\) hub blocks, indexed by
\(p=1,\ldots,L\) and by \(*\), with spectra disjoint from the
old blocks. The individual hub \(p\) sends an input arrow to
repeat position \(p\) of every old block \(i\), with fixed generic
value \(U_{i,p}\). The common hub \(*\) sends the same generic
value \(U_{i,*}\) to every repeat position in old block \(i\).
There are no reverse arrows at this test input. Arrow values for
different pairs of old blocks and hubs can be chosen independently
and to satisfy the surjectivity conditions of
Lemma~\ref{lem:pair-quotient-surjectivity}.

Figure~\ref{fig:hub-test} shows the distinction between the individual
hubs and the common hub.

\begin{figure}[tbp]
\centering
\begin{tikzpicture}[
  x=1cm,y=1cm,>=Stealth,
  every node/.style={font=\small},
  block/.style={rectangle,rounded corners=2pt,draw=black!65,fill=white,
    minimum width=20mm,minimum height=8.5mm,inner sep=4pt,align=center},
  forward/.style={draw=blue!65!black,line width=.85pt,->},
  common/.style={draw=teal!65!black,line width=.85pt,->},
  arrowlabel/.style={fill=white,inner sep=2pt}
]
  \node[font=\small\bfseries] at (6,5.9)
    {One old block $i$, with three repeat positions at $\delta=0$};
  \foreach \p/\xx in {1/1.4,2/6,3/10.6}{
    \node[block] (hub\p) at (\xx,4.7) {hub $\p$};
    \node[block] (leaf\p) at (\xx,2.35) {leaf $(i,\p)$};
    \draw[forward] (hub\p) --
      node[arrowlabel,right,pos=.5] {$U_{i,\p}$} (leaf\p);
  }
  \node[block,minimum width=25mm] (common) at (6,0) {common hub $*$};
  \draw[common] (common.west) --
    node[arrowlabel,above,sloped,pos=.55] {$U_{i,*}$} (leaf1.south);
  \draw[common] (common.north) --
    node[arrowlabel,right,pos=.5] {$U_{i,*}$} (leaf2.south);
  \draw[common] (common.east) --
    node[arrowlabel,above,sloped,pos=.55] {$U_{i,*}$} (leaf3.south);
  \node[align=center] at (6,-1.03)
    {The common hub uses the same matrix pair at every repeat position.};
  \node[align=center] at (6,-1.62)
    {Reverse arrows are zero at the base point and supply derivative directions.};
\end{tikzpicture}

\caption{The hub test at $\delta=0$, for one old block $i$ and $L=3$.
Hub $p$ sends a generic matrix pair only to position $p$; the common hub
sends the same pair to all positions. Every old block uses these same
$L+1$ hubs. At $\delta\ne0$ the leaf input may couple repeat positions.
The individual hubs eliminate off-diagonal dual entries, and the common
hub identifies the diagonal functionals.}
\label{fig:hub-test}
\end{figure}

Treat the entire deformed leaf array as one block and order it
before the hubs. Over the \(\delta\)-Laurent field this is a star
with maximal root count. Put
\[
 N=mL+L+1.
\]
Select the \(L\) branches over each old selected label and the
selected labels of all hubs, for a total of \(Na\) output rows.
The rank of their derivative matrix is at most
\(\rho_{bN}\leq Na-g\). This bound follows from the original
trace-minor identities, evaluated coefficientwise at the present
Taylor input; it does not require the specialized input to be
generic.

\newpage
\begin{lemma}[The limiting derivative columns]
\label{lem:limiting-columns}
The following derivative columns at the preceding star have finite
limits at \(\delta=0\).
\begin{enumerate}
\item A loop \(Y\otimes T\) inside old block \(i\) has limit
\[
 D_{i,h}(Y)\operatorname{tr}(\Pi_{i,h,j}T)
\]
on row \((i,h,j)\), and zero on all other rows. Hub loops give
their own loop responses.
\item For individual hub \(p\), a reverse arrow of value \(V\)
from repeat position \(r\) in old block \(i\) has limit
\begin{align*}
 &m_{i,h}(P_i,Q_p;U_{i,p},V)
      \operatorname{tr}(\Pi_{i,h,j}E_{pr})
       &&\text{on row }(i,h,j),\\
 &\mathbf 1_{p=r}\,
   m_{p,k}(Q_p,P_i;V,U_{i,p})
       &&\text{on hub row }(p,k),
\end{align*}
and zero elsewhere. Here \(E_{pr}\) is the repeat-factor matrix
unit with row \(p\) and column \(r\).
\item For the common hub, use arrow \(U_{i,*}\) and replace the
repeat-factor matrix by
\[
 \sum_{p=1}^L E_{pr}.
\]
The hub contribution is the single response
\(m_{*,k}(Q_*,P_i;V,U_{i,*})\).
\item The derivative along the whole input \(\delta\)-curve,
multiplied by \(\delta^{1-\ell}\), has limit
\(\ell\theta_{i,h,j}\) on leaf rows and zero on hub rows.
\end{enumerate}
The matrix consisting of a finite spanning collection of these
limiting columns has rank at most \(Na-g\).
\end{lemma}

\begin{proof}
Over the Laurent field, apply Lemma~\ref{lem:star-derivatives} with
the whole leaf array as one block. For loop derivatives, the
bounded leaf projection tends to
\(\pi_{i,h}\otimes\Pi_{i,h,j}\). Compression of the ordinary
first response gives \(D_{i,h}(Y)\) on the multiplicity space,
and taking trace divided by \(t\) gives the first formula.

For paired derivatives on a leaf, the Schur inverse is on the hub
at a root \(\lambda_{i,h,j}(\delta)\); that root stays away from
the hub spectrum at \(\delta=0\). For paired derivatives on a
hub, the inverse is on the leaf array at a hub root, which stays
away from every old leaf root. These inverses are regular at
\(\delta=0\). No inverse at the difference of two approaching
leaf roots occurs in this calculation.

At the reduction, the leaf input is diagonal between old blocks
and repeated inside each old block. A forward arrow at position
\(p\) followed by a reverse arrow from position \(r\) inserts
\(E_{pr}\) on the leaf repeat factor. At the hub end, the repeat
indices must match, giving \(\mathbf 1_{p=r}\). The ordinary
paired-response formula and trace divided by \(t\) now give the
second formula. Summing forward positions gives the common-hub
formula. The bounded projections and the distant-root inverses
justify all these limits coefficientwise.

The curve column lies in the span of input derivative columns over
the Laurent field by the chain rule. The expansion in
Lemma~\ref{lem:repeated-splitting} gives
\[
 \delta^{1-\ell}\frac{d}{d\delta}
       \lambda_{i,h,j}(\delta)
 =\ell\theta_{i,h,j}+O(\delta).
\]
Hub roots are constant on the curve. Thus this column also has a
finite limit. Choose bases for the finitely many loop and reverse
direction spaces and include the curve column. Every
\((Na-g+1)\)-minor of the resulting matrix vanishes over the
Laurent field; since its entries have finite limits, these minors
remain zero after specialization. This proves the last assertion.
\end{proof}

We now use the individual hubs to remove off-diagonal repeat-factor
weights and the common hub to make all diagonal weights agree.
The rank bound will then show that no element of \(\mathcal G\)
has been lost.

\begin{lemma}[The left nullspace]
\label{lem:hub-nullspace}
The left nullspace of the limiting matrix in
Lemma~\ref{lem:limiting-columns}, including its curve column,
has dimension \(g\). Its projection to the old blocks identifies
it with \(\mathcal G\): for each \(c\in\mathcal G\), its leaf
weights are \(c_{i,h}\) on every row \((i,h,j)\), independently
of \(j\), and its hub weights are the uniquely determined
compatible functionals. Consequently
\begin{equation}
 \sum_i\sum_{h\in I_i}c_{i,h}\operatorname{tr}(B_{i,h})=0
 \qquad(c\in\mathcal G).
 \label{eq:error-matrix-trace}
\end{equation}
\end{lemma}

\begin{proof}
Let \(\mu_{i,h,j}\) be leaf row weights in the left nullspace
and let \(\nu_p,\nu_*\) be the hub weights. Define
\[
 T_{i,h}=\sum_{j=1}^L\mu_{i,h,j}\Pi_{i,h,j}\in M_L(E).
\]
The loop columns imply
\[
 \sum_{h\in I_i}D_{i,h}(Y)
          \operatorname{tr}(T_{i,h}T)=0
 \qquad\text{for every }Y,T.
\]
Thus, for every matrix entry \((r,p)\), the coordinate vector
\( ((T_{i,h})_{rp})_{h\in I_i}\) defines a functional in
\(H_{P_i}^*\). The hub weights similarly belong to their own
loop annihilators.

Fix an individual hub \(p\) and a reverse position \(r\ne p\).
There is no hub response. Since
\[
 \operatorname{tr}(T_{i,h}E_{pr})=(T_{i,h})_{rp},
\]
the reverse columns say that the functional with these
\((r,p)\)-entries kills the leaf endpoint projection of the paired
map. That projection is onto \(H_{P_i}\), so all these entries
vanish. As \(p,r\) vary, every \(T_{i,h}\) is diagonal. When
\(r=p\), the corresponding diagonal-entry functional and the
hub functional \(\nu_p\) are exactly a pair in the relation of
Proposition~\ref{prop:compatible-relations}.

For the common hub, the leaf expression is
\[
 \operatorname{tr}\left(T_{i,h}\sum_pE_{pr}\right)
   =\sum_p(T_{i,h})_{rp}=(T_{i,h})_{rr}.
\]
For every \(r\), this diagonal-entry functional is therefore
paired with the same \(\nu_*\). Projection of the pair relation
onto the hub dual space is an isomorphism, so the diagonal-entry
functionals are independent of \(r\). Compatibility for the
collection of old blocks and hubs identifies them across all old
blocks. Hence
\[
 T_{i,h}=c_{i,h}I_L,
 \qquad c\in\mathcal G,
\]
and all hub weights are uniquely determined by \(c\). Since the
rank-one projections \(\Pi_{i,h,j}\) are mutually orthogonal
and sum to identity, this also gives
\(\mu_{i,h,j}=c_{i,h}\). The map from the left nullspace to
\(\mathcal G\) is consequently injective.

There are \(Na\) rows, and the limiting rank is at most \(Na-g\)
even after the curve column is included. The nullspace therefore
has dimension at least \(g\). Its injection into the
\(g\)-dimensional space \(\mathcal G\) is an isomorphism.
In particular every \(c\in\mathcal G\), not just a possible
subspace of relations, occurs. Pairing its null vector with the
curve column and cancelling the nonzero integer \(\ell\) yields
\[
 0=\sum_i\sum_{h\in I_i}\sum_{j=1}^L
       c_{i,h}\theta_{i,h,j}
  =\sum_i\sum_{h\in I_i}c_{i,h}\operatorname{tr}(B_{i,h}).
\]
If \(g=0\), the same proof gives zero nullspace and the assertion
about \(c\) is vacuous. If \(L=1\), the off-diagonal step has
no entries to remove; the common-hub step still gives precisely
the stated conclusion.
\end{proof}

\begin{proof}[Proof of Theorem~\ref{thm:cyclic-compatibility}]
Apply Lemma~\ref{lem:hub-nullspace} to arbitrary \(S_1,S_2\)
and to the auxiliary loops chosen in
Lemma~\ref{lem:repeated-splitting}. The added-loop contribution
to \eqref{eq:error-matrix-trace} is zero. Indeed, write
\(Y_i=\sum_\alpha Y_{i,\alpha}\otimes T_{i,\alpha}\). Then
\[
 \sum_{h\in I_i}c_{i,h}\operatorname{tr}(D_{i,h}(Y_i))
 =\sum_\alpha\operatorname{tr}(T_{i,\alpha})
      \sum_{h\in I_i}c_{i,h}D_{i,h}(Y_{i,\alpha})=0,
\]
because \(c_i\) annihilates the loop image. We conclude that,
for arbitrary matrix substitutions,
\begin{equation}
 \sum_{|\omega|=\ell}
 \left(\sum_i\sum_{h\in I_i}c_{i,h}e_{i,h,\omega}\right)
 \operatorname{tr}(S^\omega)=0.
 \label{eq:all-matrix-trace-tests}
\end{equation}
Auxiliary field extensions impose no restriction: this equality
has its coefficients in the original field, and extension is
injective.

We use a matrix-cycle test of the kind appearing in the constructive
trace argument of \cite[Proposition~4.1]{KlepSpenko2014}, recording the
rotational multiplicity explicitly. Fix a word
\(\omega_0=c_0c_1\cdots c_{\ell-1}\), with letters
\(c_j\in\{1,2\}\), and take \(L=\ell\). On vertices indexed
modulo \(\ell\), set
\[
 S_c=\sum_{j:\,c_j=c}E_{j,j+1}\qquad(c=1,2).
\]
For a word \(\omega\) of length \(\ell\), its matrix trace
counts the starting vertices at which one reads \(\omega\)
around this directed cycle. It is zero unless \(\omega\) lies
in the cyclic orbit \(\mathcal O\) of \(\omega_0\); for every
distinct word in that orbit its value is the same positive integer
\[
 \frac{\ell}{|\mathcal O|}.
\]
This is the rotational stabilizer size and includes periodic and
constant-letter words. Substitution into
\eqref{eq:all-matrix-trace-tests}, followed by cancellation of that
integer in characteristic zero, gives
\eqref{eq:cyclic-error-relation}.
\end{proof}

The leading errors now satisfy every cyclic relation imposed by
\(\mathcal G\). The colored correction graph will make these
exactly the annihilators of its order-by-order correction map.
